\section{Porteurs du projet}

% ----------------------------------------------------------------
% 📌 À quoi sert cette section ?
% ----------------------------------------------------------------
% Cette partie sert à identifier formellement les étudiants responsables du projet.
% Elle permet de les associer clairement à la production, pour le suivi et la validation.

% ----------------------------------------------------------------
% 🧭 Comment la remplir ?
% ----------------------------------------------------------------
% ➤ Indiquez les informations suivantes pour chaque porteur :
%     - Prénom NOM
%     - Promotion (ex: M1, M2)
%     - Campus SUPINFO
%     - Adresse e-mail institutionnelle (prenom.nom@supinfo.com)

% ➤ Si vous êtes plusieurs, vous pouvez organiser l’information sous forme de liste ou de tableau.

% ----------------------------------------------------------------
% 🧾 Exemple sous forme de liste :
% ----------------------------------------------------------------
% \begin{itemize}
%   \item Rémi COVIZZI – M2 – Lyon – remi.covizzi@supinfo.com
%   \item Sophie MARTIN – M1 – Paris – sophie.martin@supinfo.com
% \end{itemize}

% ----------------------------------------------------------------
% 📊 Exemple sous forme de tableau :
% ----------------------------------------------------------------
% \begin{table}[H]
% \centering
% \begin{tabular}{|l|c|c|l|}
% \hline
% Nom & Promotion & Campus & Email \\
% \hline
% Rémi COVIZZI & M2 & Lyon & remi.covizzi@supinfo.com \\
% Sophie MARTIN & M1 & Paris & sophie.martin@supinfo.com \\
% \hline
% \end{tabular}
% \caption{Étudiants porteurs du projet}
% \end{table}

% ----------------------------------------------------------------
% ✍️ À vous de jouer : complétez ci-dessous
% ----------------------------------------------------------------

% (Modèle recommandé)
\begin{itemize}
  \item [À compléter]
\end{itemize}
